\section{Stage 230: Relaxed stationary barrier front end}
\label{sec:app-part08:relaxed-stationary-barrier-front-end}
\label{stage:230}
\stagefield{Verification}{SymPy audit: \StageFile{scripts/moving_throat_pde_stage230_relaxed_stationary_barrier_compiler_from_one_port_short_range_leakage_uv_and_compensated_source_packets_sympy_audit.py}.  Mathematica audit: none yet.}


Stage 230 assembles the first stationary lowered barrier from the one-port short-range baseline and the three relaxed packets.  The one-port static bundle is controlled by
\begin{equation}
\label{eq:app-part08-stage230-delta-q-p}
\Delta=\Omega_U^2\Omega_W^2-R_{\rm mix}^{2},
\qquad
Q=G_U^2\Omega_W^2+2G_UG_WR_{\rm mix}+G_W^2\Omega_U^2,
\end{equation}
\begin{equation}
\label{eq:app-part08-stage230-p-d0}
P=\Omega_U^2G_W+R_{\rm mix}G_U,
\qquad
D_0=K_*-\frac{Q}{\Delta}.
\end{equation}
With primitive short-range source amplitudes \(\beta_Q,\beta_U,\beta_W\), the product-family coefficients are
\begin{equation}
\label{eq:app-part08-stage230-c6c4c2}
\mathcal C_6=\chi_{qq}\beta_Q^2,
\qquad
\mathcal C_4=\chi_{qU}\beta_Q\beta_U+\chi_{qW}\beta_Q\beta_W,
\end{equation}
\begin{equation}
\label{eq:app-part08-stage230-c2}
\mathcal C_2=\chi_{UU}\beta_U^2+2\chi_{UW}\beta_U\beta_W+\chi_{WW}\beta_W^2.
\end{equation}
The collected baseline is
\begin{equation}
\label{eq:app-part08-stage230-vshort}
V_{\rm short}^{(1p)}(r)
=
\frac1r\left(1+\frac12e^{-2\kappa r}\right)
-
\frac{A_6}{r^6}
-
A_4\frac{e^{-2\kappa r}}{r^4}
-
A_2\frac{e^{-4\kappa r}}{r^2},
\end{equation}
where
\begin{equation}
\label{eq:app-part08-stage230-a6a4a2}
A_6=3\alpha_6+\frac12\mathcal C_6,
\qquad
A_4=\mathcal C_4,
\qquad
A_2=\alpha_2+\frac12\mathcal C_2.
\end{equation}
This baseline inherits the same short-range firewall from Stages 202--226.

The imported relaxed scalars are
\begin{equation}
\label{eq:app-part08-stage230-lag-l}
\mathfrak L(r):=\Lambda(r)\varrho(r),
\qquad
S_{\rm leak}=\frac{8\sqrt2\eta_{\rm leak}\mu_w\rho_0}{\pi^{5/2}\lambda^3}\mathfrak L,
\end{equation}
\begin{equation}
\label{eq:app-part08-stage230-work}
\mathcal W_w^{\rm sess}
=
\frac{512\eta_{\rm leak}^{2}\mu_wq\rho_0}{\pi^4\lambda^2}\mathfrak L^2,
\qquad
\Delta E_{UV}=\eta_{UV}\mathcal D_{UV},
\end{equation}
\begin{equation}
\label{eq:app-part08-stage230-source-lowering}
\mathcal M_\sigma(r)=\xi_R\bigl[\mathcal R_\infty-\mathcal R(r)\bigr],
\qquad
\mathcal R_\infty=\frac{(2/\pi-\mathfrak r_{F1})^2}{1+\mathfrak r_{F1}^2}.
\end{equation}
The stationary compiler is \cref{eq:app-part08-intro-veff230}; the lowering identity is \cref{eq:app-part08-main-lowering-identity}.

The Session-I benchmark is useful as an audit point, not as a general theorem.  At the strongest softening point, the Stage-230 decomposition reads
\begin{equation}
\label{eq:app-part08-stage230-benchmark-decomposition}
3.74163698
-
0.26971918\times0.31069599
-
1.51632107
-
0.21064278
-
0.18386120
=
1.74701126.
\end{equation}
This shows that, once the one-port baseline is fixed, the recorded stationary lowering can be decomposed into work, \(U/V\) drain, compensated source response, and one remaining leakage-to-barrier embedding coefficient.  It does not change the short-range kernel class.
